\section{Threat Model and Assumptions}
\label{sec:threat}

We consider three adversaries.
\textbf{A1}, a persuasive vendor, controls everything a resource server returns, including prices, payees
and descriptions, and aims to make the agent pay more, more often, or to someone else.
\textbf{A2} has arbitrary code execution in the agent's process, the \emph{orchestrator}: it holds the relay
key, relays everything between the signer and the user's console, and may submit any request to the policy
contract and any query to the signer, in any order. A2 subsumes A1, and our claims about the agent are made
against A2. \textbf{A3}, a dishonest vendor, collects a payment after its hold has ended, or collects in
time and delivers nothing or something worthless. The implemented system addresses late funding;
non-delivery is the subject of Section~\ref{sec:delivery}.

\paragraph{What the adversaries cannot do.} Break Keccak-256 or ECDSA over secp256k1; obtain the user's
wallet key; rewrite the canonical chain. The signer's RPC endpoint may lie, and blocks may be dropped from
the chain it reads; for payment authorization after a stable, authenticated setup this costs liveness and
cannot widen the loss bound (Section~\ref{sec:escrow}). The leakage theorem makes a stronger assumption
about dropped blocks, stated in it (Section~\ref{sec:leakthm}).

\paragraph{Trusted components.} The loss bound assumes an uncompromised signer: the payer is an ordinary
key, so a compromised signer can sign a plain transfer and take the funded balance. It also trusts the
ROFL application's administrator, who can change which enclave measurements receive the application's
keys. The user's console and the user's own RPC endpoint are trusted; the console is built with the vault
address and the signer's enrollment address fixed.

\paragraph{Inco.} We separate what we assume of Inco from what we test, since a verified signature
establishes only that a threshold of covalidators endorsed a plaintext. We assume:
\begin{enumerate}[(I1)]
  \item \emph{Correct evaluation:} each handle an operation returns evaluates to that operation applied to
        the plaintexts of its inputs.
  \item \emph{Ordered evaluation and binding:} the operation graph is evaluated in the order of the
        canonical chain's transactions and events, so that the budget handle a request reads is the one
        the previous transaction wrote, and the decision handle the vault stores for $(g, s)$ is the one
        that request computed.
  \item \emph{Unforgeable attestation:} a threshold signature verifies for $(h, v)$ only if handle $h$
        evaluates to $v$.
  \item \emph{Access control:} only the vault can use the budget and cap handles as operands.
  \item \emph{Confidentiality:} no plaintext of a budget or cap handle reaches anyone; only the decisions the
        vault reveals are decrypted.
\end{enumerate}
The loss bound needs I1 to I4; the leakage theorem needs I5 as well. What this paper tests is the
contract logic \emph{given} these properties, using Inco's test double, and the live behaviour of the deployed
service in our runs; neither establishes the properties themselves. If I1 or I3 fails, the public call count,
the allowlist and the funded balance still bound the loss.
